\documentclass[11pt]{article}
\usepackage{fontspec}
\setmainfont{Times New Roman}
\setsansfont{Arial}
\setmonofont{Consolas}
\usepackage{amsmath,amssymb,mathtools}
\usepackage{geometry}
\usepackage{microtype}
\geometry{a4paper,margin=2.35cm}
\setlength{\parindent}{1.25em}
\setlength{\parskip}{0pt}
\emergencystretch=25em
\sloppy
\hbadness=10000
\vbadness=10000
\hfuzz=2pt
\vfuzz=10pt
\raggedbottom

\begin{document}

% Unit: Paper26 v001. Direct Ukrainian translation from the German final-audited source slice.
\section*{26. Абстрактна побудова теорії ідеалів в алгебраїчному числовому полі}

\begin{center}
J. Ber. d. DMV 33 (1924), с. 102
\end{center}

4. E. Noether, Геттінген: Абстрактна побудова теорії ідеалів в алгебраїчному числовому полі.

Було вказано абстрактні властивості, які цілком еквівалентні теорії ідеалів в алгебраїчному числовому полі і, отже, характеризують усі кільця, що мають таку теорію ідеалів. Це кільця з одиницею і без дільників нуля, у яких виконується умова подвійного ланцюга: кожен ланцюг ідеалів за діленням обривається після скінченної кількості кроків, і так само кожен ланцюг кратних за модулем кожного ідеалу, відмінного від нульового ідеалу; крім того, ці кільця є інтегрально замкненими відносно свого поля часток. Якщо відкинути останнє припущення, то отримуємо твердження про теорію ідеалів у скінченному порядку, які частково вже трапляються без доведення у Dedekind (3-тє вид.).

Докладний виклад з'явиться в Mathematische Annalen.

\end{document}
